% Population-level Jansen--Rit diagram; synaptic filters sit at receiving boundaries.
\begingroup%
\definecolor{jrPyramidal}{HTML}{4477AA}\definecolor{jrExcitatory}{HTML}{009E73}\definecolor{jrInhibitory}{HTML}{D55E00}%
\tikzset{>=Stealth,
 thesis figure/.style={font=\rmfamily\small},
 thesis heading/.style={font=\rmfamily\small\bfseries},
 thesis secondary/.style={font=\rmfamily\footnotesize},
 flow/.style={->,line width=1.1pt,rounded corners=4pt},
 card/.style={rounded corners=7pt,draw=black!25,fill=white,align=center},
 pics/pyramidal/.style={code={
 \draw[line width=.75pt] (0,.14)--(0,.85) (0,.5)--(-.26,.73) (0,.57)--(.25,.8)
 (0,.85)--(-.12,1.05) (0,.85)--(.12,1.05)
 (0,-.12)--(0,-.65) (0,-.43)--(-.22,-.65) (0,-.43)--(.22,-.65)
 (-.1,-.04)--(-.45,.18) (-.3,.09)--(-.46,.4)
 (.1,-.04)--(.45,.18) (.3,.09)--(.46,.4);
 \path[draw,fill=white,line width=.9pt] (-.17,-.13)--(.17,-.13)--(0,.18)--cycle;
 }},
 pics/stellate/.style={code={
 \foreach \a in {0,60,120,180,240,300} {
% TikZ angle arithmetic uses a minus sign.
 \draw[line width=.75pt] (\a:.12)--(\a:.47)--(\a+15:.66)
 (\a:.47)--(\a-15:.66);} % latex-lint:ignore=TYPO-05
 \draw[fill=white,line width=.9pt] (0,0) circle(.15);
 }},
 pics/inhibitory/.style={code={
 \foreach \a in {30,120,210,300} {
% TikZ angle arithmetic uses a minus sign.
 \draw[line width=.75pt] (\a:.13)--(\a:.4)--(\a+20:.6)
 (\a:.4)--(\a-20:.6);} % latex-lint:ignore=TYPO-05
 \draw[line width=.75pt] (0,-.15)..controls(.1,-.5)and(.6,-.4)..(.5,-.75);
 \draw[fill=white,line width=.9pt] (0,0) circle(.17);
 }}}
% Repeated glyphs denote an ensemble, not a literal number of cells.
\newcommand{\ensemble}[3]{\begin{scope}[shift={#1},draw=#2]
% The cell ensemble uses a negative TikZ coordinate offset.
 \foreach \x/\y in {-1.5/0,0/0,1.5/0}{ % latex-lint:ignore=TYPO-05
 \pic[scale=.43] at (\x,\y) {#3};}
 \end{scope}}%

\definecolor{potential}{HTML}{0a0a0a}%
\definecolor{rate}{HTML}{AA3377}%
\tikzset{
 potential/.style={flow,draw=potential},
 rate/.style={flow,draw=rate},
 op/.style={draw=black!75,fill=white,rectangle,rounded corners=1pt,minimum width=1.35cm,minimum height=.72cm,align=center,line width=.9pt},
 mul/.style={op,minimum width=1.0cm,minimum height=.65cm},
 sum/.style={circle,draw=black!75,fill=white,minimum size=.55cm,inner sep=0pt},
 smalllabel/.style={thesis secondary,fill=black!1,inner sep=2pt}}%
\newcommand{\signallegend}[1]{
% The legend uses a negative TikZ coordinate offset.
 \draw[potential] (#1) ++(-4.7,0)--++(1,0) node[right,black,thesis secondary] {Mean potential (mV)}; % latex-lint:ignore=TYPO-05
 \draw[rate] (#1) ++(1.5,0)--++(1,0) node[right,black,thesis secondary] {Mean firing rate ($\mathrm{s}^{-1}$)};}%

\begin{tikzpicture}[thesis figure,y=1.15cm]

\node[card,fill=white,minimum width=16.8cm,minimum height=11.5cm] at (7.8,0) {};
\node[thesis heading] at (7.8,5.45) {One Jansen--Rit node};
\node[card,draw=jrExcitatory!55,fill=jrExcitatory!8,minimum width=5cm,minimum height=2.9cm] at (5.5,3.3) {};
\node[card,draw=jrInhibitory!55,fill=jrInhibitory!8,minimum width=5cm,minimum height=2.9cm] at (5.5,-3.3) {};
\node[card,draw=jrPyramidal!55,fill=jrPyramidal!8,minimum width=8.7cm,minimum height=3cm] at (8.65,0) {};
\node[thesis heading,text=jrExcitatory] at (5.5,4.15) {Excitatory interneurons};
\node[thesis heading,text=jrInhibitory] at (5.5,-4.15) {Inhibitory interneurons};
\node[thesis heading,text=jrPyramidal] at (8.6,.9) {Pyramidal population};
\begin{scope}[shift={(4.6,3.6)},scale=.55]\ensemble{(0,0)}{jrExcitatory}{stellate}\end{scope}
\begin{scope}[shift={(4.6,-3.65)},scale=.55]\ensemble{(0,0)}{jrInhibitory}{inhibitory}\end{scope}
\begin{scope}[shift={(5.35,.65)},scale=.55]\ensemble{(0,0)}{jrPyramidal}{pyramidal}\end{scope}
\node[op] (ef) at (3,3) {$h_{\mathrm e}$};
\node[op] (if) at (3,-3) {$h_{\mathrm e}$};
\node[op] (es) at (6.8,3) {$S$};
\node[op] (is) at (6.8,-3) {$S$};
\node[mul] (c1) at (.6,1.5) {$c_1$};
\node[mul] (c3) at (.6,-1.5) {$c_3$};
\node[mul] (c2) at (10,3) {$c_2$};
\node[mul] (c4) at (10,-3) {$c_4$};
\node[sum] (drive) at (11.75,3) {$\Sigma$};
\node[op] (he) at (13,.7) {$h_{\mathrm e}$};
\node[op] (hi) at (13,-.7) {$h_{\mathrm i}$};
\node[sum] (bal) at (10.8,0) {$\Sigma$};
\node[sum] (stim) at (8.8,0) {$\Sigma$};
\node[op] (ps) at (5.3,0) {$S$};
\draw[rate,-] (ps.west)--(.6,0) node[pos=.6,below,smalllabel,text=rate] {$r_{\mathrm p}$};
\draw[rate] (.6,0)--(c1.south);
\draw[rate] (c1.north)--(.6,3)--(ef.west);
\draw[rate] (.6,0)--(c3.north);
\draw[rate] (c3.south)--(.6,-3)--(if.west);
\fill[fill=rate] (.6,0) circle(2pt);

\draw[potential] (ef.east)--(es.west) node[midway,below,smalllabel,fill=jrExcitatory!8,text=potential] {$c_1x_{i,1}$};
\draw[potential] (if.east)--(is.west) node[midway,above,smalllabel,fill=jrInhibitory!8,text=potential] {$c_3x_{i,1}$};
\draw[rate] (es.east)--(c2.west) node[pos=.75,above,smalllabel,text=rate] {$r_{\mathrm e}$};
\draw[rate] (is.east)--(c4.west) node[pos=.75,below,smalllabel,text=rate] {$r_{\mathrm i}$};
\draw[rate] (c2.east)--(drive.west);
\draw[rate] (11.75,4) node[above,black] {$I+I_{\mathrm{coup},i}+\zeta_i$} --(drive.north);
\draw[rate] (drive.east)--(15,3)|-(he.east);
\draw[rate] (c4.east)--(15,-3)|-(hi.east);
\draw[potential] (he.west)--(10.8,.7) node[pos=.25,above,smalllabel,fill=jrPyramidal!8,text=potential] {$x_{i,2}$}
    --(bal.north) node[pos=.99,right,black] {$+$};
\draw[potential] (hi.west)--(10.8,-.7) node[pos=.25,below,smalllabel,fill=jrPyramidal!8,text=potential] {$x_{i,3}$}
    --(bal.south) node[pos=.99,right,black] {$-$};
\draw[potential] (bal.west)--(stim.east);
\draw[potential] (9.8,0)--(9.8,-1.8) node[below,black] {$y_{\mathrm{reg},i}$};
\draw[potential] (8.8,-1.8) node[below,black] {$s_i$} --(stim.south) node[pos=.95,right,black] {$+$};
\draw[potential] (stim.west)--(ps.east) node[midway,below,smalllabel,fill=jrPyramidal!8,text=potential] {$x_{i,2}-x_{i,3}+s_i$};
\fill[fill=potential!100] (9.8,0) circle(2pt);
\signallegend{7.2,-5.9}
\end{tikzpicture}%
\endgroup%
